% Generated draft: requires mathematical and source review.
\documentclass[11pt]{article}
\usepackage[margin=25mm]{geometry}
\usepackage{fontspec}
\setmainfont{Latin Modern Roman}
\usepackage{amsmath,amssymb,mathtools,unicode-math}
\setmathfont{Latin Modern Math}
\usepackage{enumitem,needspace,xurl,hyperref}
\setlength{\parindent}{0pt}
\setlength{\parskip}{5pt}
\setlength{\emergencystretch}{3em}
\newcommand{\N}{\mathbb{N}}
\newcommand{\Z}{\mathbb{Z}}
\newcommand{\Q}{\mathbb{Q}}
\newcommand{\R}{\mathbb{R}}
\newcommand{\C}{\mathbb{C}}
\newcommand{\F}{\mathbb{F}}
\newcommand{\A}{\mathbb{A}}
\newcommand{\PP}{\mathbb{P}}
\newcommand{\E}{\mathbb{E}}
\newcommand{\eps}{\varepsilon}
\newcommand{\dd}{\,\mathrm{d}}
\begin{document}
\section*{UnsolvedMath ID: 30005942}
\textbf{Source problem code:} catalogue entry
\subsection{Definitions and mathematical statement}
Let \(S\) denote a surface in the class intended by the source, let \(\omega_S\) denote its canonical line bundle, and write
\[
K_S:=\operatorname{Tot}(\omega_S).
\]
The source describes \(K_S\) as a Calabi--Yau threefold. Write \(\operatorname{BPS}_{\beta_S}^{S}(q)\) for the generalized Block--Göttsche polynomial attached in the source to \(S\) and the class label \(\beta_S\), with polynomial variable \(q\).

\textbf{Conjectural target.} For each surface \(S\) and class \(\beta_S\) permitted by the source's construction, do the polynomials
\[
\operatorname{BPS}_{\beta_S}^{S}(q)
\]
admit an interpretation in refined Donaldson--Thomas theory of \(K_S\)? The supplied literature assessment describes the intended interpretation as an identification with refined BPS invariants of \(K_S\).

\textbf{Unresolved scope and definitions.} The supplied record does not specify the permitted surfaces, the admissible classes \(\beta_S\), the construction or normalization of these polynomials, or the precise refined Donaldson--Thomas invariant and variable convention to be compared with them. Thus the preceding quantification refers to an unspecified source-defined range; it cannot be replaced here by an explicit universal quantifier over a determined class of surfaces and classes. A precise invariant equality requires those missing definitions.

\textbf{Reported special case and qualification.} The original excerpt says that degeneration arguments prove ``this conjecture'' when \(S\) is a del Pezzo surface, described there as the case \(n\leq 8\). It then states a theorem attributed to Arg\"uz--Bousseau:
\[
\operatorname{BPS}_{\beta_S}^{S}(-1)=W_{\beta_S}
\qquad\text{for every }n\leq 8.
\]
Here \(n\) and \(W_{\beta_S}\) are retained as source notation: their definitions, admissible classes, and any real-geometric or enumerative choices are absent from the supplied excerpt. The relationship between this specialization identity and the full refined Donaldson--Thomas interpretation is not specified. In particular, the excerpt's claim to prove ``this conjecture'' in the del Pezzo case must be reconciled with the supplied dated assessment that the proposed identification remains conjectural and only related special evidence is established.
\subsection{Short English statement}
The question asks whether generalized Block--Göttsche polynomials associated with a surface \(S\) can be identified with suitable refined Donaldson--Thomas BPS invariants of the total space of its canonical line bundle. The permitted surfaces, classes, and precise invariants are unspecified. The source also reports an Arg\"uz--Bousseau identity at \(q=-1\) for \(n\leq 8\), but its connection to the full conjecture requires clarification.
\subsection{Sources}
\begin{itemize}
\item Local UnsolvedMath registry, permanent ID \(30005942\), source problem code \texttt{catalogue entry}, title: \emph{Refined Donaldson--Thomas Meaning of Block--Göttsche Polynomials}. The supplied record assigns proposed year 2024 and names no proposer. Its status is recorded as open, and its statement status as \texttt{corrected\_verified}; the accompanying claim of verification against the official report is supplied metadata, not an independent verification in this draft.
\item \emph{Real and Logarithmic Enumerative Geometry} (2024), Oberwolfach Reports, Report 23/2024, pp. 1285--1338, DOI 10.4171/owr/2024/23. Supplied report reference and literature source: \url{https://doi.org/10.4171/owr/2024/23}. The supplied original excerpt attributes the displayed specialization theorem to Arg\"uz--Bousseau.
\item \emph{BPS polynomials and Welschinger invariants}. Supplied literature source: \url{https://arxiv.org/abs/2506.02770}.
\item Supplied literature assessment, checked on \textbf{2026-08-22}: the proposed identification with refined BPS invariants of the canonical-line-bundle Calabi--Yau remains conjectural, with only related special evidence established. This is supplied dated curation derived from the Oberwolfach material, not independently verified here; time-sensitive status and cited claims require independent checking.
\end{itemize}
\end{document}
